A modularidade é um atributo importante para a qualidade estrutural de sistemas de software, pois está relacionada à 
forma como responsabilidades são distribuídas, dependências são estabelecidas e componentes podem ser compreendidos, 
modificados e evoluídos. Em sistemas orientados a objetos, decisões arquiteturais podem influenciar diretamente características 
como acoplamento, coesão e complexidade estrutural. Nesse contexto, estilos arquiteturais como a arquitetura em camadas e a 
arquitetura hexagonal propõem diferentes formas de organizar os elementos do sistema, podendo produzir efeitos distintos sobre 
sua modularidade.

Diante disso, este trabalho tem como objetivo avaliar a modularidade de sistemas orientados a objetos por meio da aplicação de 
métricas estruturais em diferentes estilos arquiteturais. Para isso, será conduzido um estudo comparativo envolvendo 
dois sistemas desenvolvidos em contexto acadêmico: o SISO, estruturado segundo a arquitetura em camadas, e o Roomie, estruturado 
segundo a arquitetura hexagonal. A análise será realizada a partir de métricas extraídas por análise estática do código-fonte, 
considerando indicadores relacionados a acoplamento, coesão e complexidade estrutural.

Espera-se que este trabalho contribua para a compreensão de como diferentes decisões arquiteturais podem se refletir em características 
estruturais mensuráveis de sistemas orientados a objetos. Além disso, busca-se fornecer evidências quantitativas que apoiem a análise da 
modularidade e auxiliem a adoção de práticas arquiteturais mais fundamentadas na Engenharia de Software.

\palavraschave{Modularidade; Arquitetura de Software; Métricas de Software; Sistemas Orientados a Objetos.}